% TODO make hw stylesheet
\documentclass[10pt]{scrartcl}
\usepackage[a4paper, total={6in, 8in}]{geometry}
\usepackage[usenames,dvipsnames,svgnames]{xcolor}
\usepackage[shortlabels]{enumitem}
\usepackage[framemethod=TikZ]{mdframed}
\usepackage{amsmath,amssymb,amsthm}
\usepackage[colorlinks]{hyperref}
\usepackage{multicol}
\usepackage{tabularx}

\hypersetup{
	colorlinks=true,
	linkcolor=blue,
	filecolor=magenta,      
	urlcolor=magenta,
	pdftitle={MAT 322 HW}
}

\usepackage{microtype}
\usepackage{mathtools}
\usepackage[headsepline]{scrlayer-scrpage}
\usepackage{thmtools}
\usepackage[textsize=footnotesize]{todonotes}

\mdfdefinestyle{mdbluebox}{roundcorner=10pt,innerbottommargin=9pt,
	linecolor=blue,backgroundcolor=TealBlue!5,}
\declaretheoremstyle[headfont=\sffamily\bfseries\color{MidnightBlue},
mdframed={style=mdbluebox},]{thmbluebox}

\mdfdefinestyle{mdbloobox}{frametitlefont=\bfseries,innerbottommargin=8pt,
	nobreak=true,backgroundcolor=TealBlue!5,linecolor=blue,}
\declaretheoremstyle[headfont=\bfseries\color{MidnightBlue},
mdframed={style=mdbloobox},headpunct={\\[3pt]},postheadspace=0pt,]{thmbloobox}

\mdfdefinestyle{mdredbox}{frametitlefont=\bfseries,innerbottommargin=8pt,
	nobreak=true,backgroundcolor=Salmon!5,linecolor=RawSienna,}
\declaretheoremstyle[headfont=\bfseries\color{RawSienna},
mdframed={style=mdredbox},headpunct={\\[3pt]},postheadspace=0pt,]{thmredbox}

\mdfdefinestyle{mdgreenbox}{linecolor=ForestGreen,backgroundcolor=ForestGreen!5,
	linewidth=2pt,rightline=false,leftline=true,topline=false,bottomline=false,}
\declaretheoremstyle[headfont=\bfseries\sffamily\color{ForestGreen!70!black},
mdframed={style=mdgreenbox},headpunct={ --- },]{thmgreenbox}

\mdfdefinestyle{mdorangebox}{linecolor=RedOrange,backgroundcolor=RedOrange!5,
	linewidth=2pt,rightline=false,leftline=true,topline=false,bottomline=false,}
\declaretheoremstyle[headfont=\bfseries\sffamily\color{RedOrange!70!black},
mdframed={style=mdorangebox},headpunct={ --- },]{thmorangebox}

\mdfdefinestyle{mdblackbox}{linecolor=black,backgroundcolor=RedViolet!5!gray!5,
	linewidth=3pt,rightline=false,leftline=true,topline=false,bottomline=false,}
\declaretheoremstyle[mdframed={style=mdblackbox}]{thmblackbox}

\declaretheoremstyle[spaceabove=6pt,spacebelow=6pt]{basehead}

\declaretheorem[style=basehead,name=Problem]{problem}
\declaretheorem[style=thmredbox,name=Problem,numbered=no]{problem*}
\declaretheorem[style=thmbloobox,name=Setup]{setup} % for config geo
\declaretheorem[style=thmredbox,name=Example,sibling=problem]{example}
\declaretheorem[style=thmredbox,name=$\bigstar$ Problem,sibling=problem]{niceprob}
\declaretheorem[style=thmbluebox,name=Theorem,sibling=problem]{theorem}
\declaretheorem[style=thmbluebox,name=Corollary,sibling=theorem]{corollary}
\declaretheorem[style=thmbluebox,name=Proposition,sibling=theorem]{prop}
\declaretheorem[style=thmbluebox,name=Lemma,numberwithin=problem]{lemma}
\declaretheorem[style=thmbluebox,name=Theorem,numbered=no]{theorem*}
\declaretheorem[style=thmbluebox,name=Lemma,numbered=no]{lemma*}
\declaretheorem[style=thmgreenbox,name=Claim,numberwithin=problem]{claim}
\declaretheorem[style=thmgreenbox,name=Claim,numbered=no]{claim*}
\declaretheorem[style=thmblackbox,name=Remark,numberwithin=problem]{remark}
\declaretheorem[style=thmblackbox,name=Remark,numbered=no]{remark*}
\declaretheorem[style=thmgreenbox,name=Definition,sibling=theorem]{definition}
\declaretheorem[style=thmgreenbox,name=Definition,numbered=no]{definition*}
\declaretheorem[style=thmorangebox,name=Warning,sibling=theorem]{warning}
\declaretheorem[style=thmorangebox,name=Warning,numbered=no]{warning*}
\declaretheorem[style=thmorangebox,name=Note,numbered=no]{note}
\declaretheorem[style=thmblackbox,name=Exercise,sibling=problem]{exercise}
\declaretheorem[style=thmblackbox,name=Exercise,numbered=no]{exercise*}
\declaretheorem[style=thmblackbox,name=Question,sibling=problem]{question}
\declaretheorem[style=thmblackbox,name=Question,numbered=no]{question*}

\addtolength{\textheight}{3.14cm}
\providecommand{\ol}{\overline}
\providecommand{\ul}{\underline}
\providecommand{\eps}{\varepsilon}
\providecommand{\half}{\frac{1}{2}}
\providecommand{\dang}{\measuredangle} %% Directed angle
\providecommand{\CC}{\mathbb C}
\providecommand{\FF}{\mathbb F}
\providecommand{\NN}{\mathbb N}
\providecommand{\QQ}{\mathbb Q}
\providecommand{\RR}{\mathbb R}
\providecommand{\ZZ}{\mathbb Z}
\providecommand{\dg}{^\circ}
\providecommand{\ii}{\item}
\providecommand{\setdiff}{\smallsetminus}
\providecommand{\alert}[1]{{\sffamily\textbf{\textcolor{blue}{#1}}}}
\providecommand{\inv}{^{-1}}
\providecommand{\mb}[1]{\mathbf{#1}}
\providecommand{\dif}{\;d}

\reversemarginpar

\newenvironment{soln}{\begin{proof}[Solution]}{\end{proof}}
\newenvironment{walkthrough}{\noindent \textbf{Walkthrough.}}{\medskip}
\newlist{walk}{enumerate}{3}
\setlist[walk]{label=\bfseries (\alph*)}

\providecommand{\pow}{\operatorname{Pow}}
\providecommand{\vocab}[1]{{\textbf{\textcolor{ForestGreen}{#1}}}}
\providecommand{\lcm}{\operatorname{lcm}}
\providecommand{\abs}[1]{\left|#1\right|}

\usepackage{asymptote}
\begin{asydef}
	defaultpen(fontsize(10pt));
	unitsize(1cm);
	size(8cm); // set a reasonable default
	usepackage("amsmath");
	usepackage("amssymb");
	settings.tex="pdflatex";
	settings.outformat="pdf";
	// Replacement for olympiad+cse5 which is not standard
	import geometry;
	// recalibrate fill and filldraw for conics
	void filldraw(picture pic = currentpicture, conic g, pen fillpen=defaultpen, pen drawpen=defaultpen)
	{ filldraw(pic, (path) g, fillpen, drawpen); }
	void fill(picture pic = currentpicture, conic g, pen p=defaultpen)
	{ filldraw(pic, (path) g, p); }
	// some geometry
	pair foot(pair P, pair A, pair B) { return foot(triangle(A,B,P).VC); }
	pair orthocenter(pair A, pair B, pair C) { return orthocentercenter(A,B,C); }
	pair centroid(pair A, pair B, pair C) { return (A+B+C)/3; }
	// cse5 abbreviations
	path CP(pair P, pair A) { return circle(P, abs(A-P)); }
	path CR(pair P, real r) { return circle(P, r); }
	pair IP(path p, path q) { return intersectionpoints(p,q)[0]; }
	pair OP(path p, path q) { return intersectionpoints(p,q)[1]; }
	path Line(pair A, pair B, real a=0.6, real b=a) { return (a*(A-B)+A)--(b*(B-A)+B); }
	// cse5 more useful functions
	picture CC() {
		picture p=rotate(0)*currentpicture;
		currentpicture.erase();
		return p;
	}
	pair MP(Label s, pair A, pair B = plain.S, pen p = defaultpen) {
		Label L = s;
		L.s = "$"+s.s+"$";
		label(L, A, B, p);
		return A;
	}
	pair Drawing(Label s = "", pair A, pair B = plain.S, pen p = defaultpen) {
		dot(MP(s, A, B, p), p);
		return A;
	}
	path Drawing(path g, pen p = defaultpen, arrowbar ar = None) {
		draw(g, p, ar);
		return g;
	}
\end{asydef}

\usepackage{mathrsfs}

\ihead{\footnotesize MAT 322 HW08}
\ohead{\footnotesize Last compiled \today}

\title{\vspace{-2em}MAT 322 HW08}
\author{\large Aditya Pahuja}
\date{\large Due 04/15}
\begin{document}
\maketitle

\begin{problem}
	Let $U$ be an open set in $\RR^k$ and $f\colon U\to\RR$
	be a $C^\infty$ function. Show that the graph of $f$
	is a $k$-dimensional $C^\infty$ submanifold of $\RR^{k+1}$.
\end{problem}

\begin{soln}
	Consider the projection map $\pi\colon \RR^{k+1}\to\RR^k$
	given by $\pi(x_1, \dots, x_k,x_{k+1}) = (x_1,\dots,x_k)$.
	The derivative of this map is the $k\times(k+1)$ matrix
	\[(\mathrm d\pi)_{x} = \begin{pmatrix}
		1 & 0 & \dots & 0 & 0\\
		0 & 1 & \dots & 0 & 0\\
		\vdots & \vdots & \ddots & \vdots& \vdots\\
		0&0&\dots & 1&0
	\end{pmatrix}\]
	so we can clearly see that $\pi$ is $C^\infty$ and has rank $k$ everywhere.
	Letting $\Gamma=\{(x,f(x))\in\RR^k\times\RR : x\in U\}$ be the graph of $f$,
	we see that if $p\in \Gamma$, then $f(U)$ is a neighborhood of $p$ such that
	$\pi(f(U)) = U$ is open in $\RR^k$; noting the properties of $\pi$
	mentioned above, this is enough to verify
	that $\Gamma$ is a $k$-dimensional $C^\infty$ manifold.
\end{soln}

\begin{problem}
	Let $c > a > 0$ be real. Consider the $C^\infty$ function
	$f\colon\{(x,y,z)\in\RR^3 : (x,y)\ne0\}\to\RR$ defined by
	\[f(x,y,z) = (c - \sqrt{x^2+y^2})^2 + z^2.\]
	Show that $f\inv(a^2)$ is a $C^\infty$ submanifold of $\RR^3$
	and describe it.
\end{problem}

\begin{soln}
	We can easily verify that $f$ is a $C^\infty$ function
	via some tedious computation.
	
	Clearly $f\inv(a^2) = g\inv(0)$ where $g(x,y,z) = f(x,y,z)-a^2$,
	so we can say that $f\inv(a^2)$ is a $C^\infty$ submanifold of $\RR^n$ if
	$(\mathrm dg)_x$ is surjective for all $x\in f\inv(0)$,
	or equivalently $(\mathrm df)_x$ is surjective for all $x\in f\inv(a^2)$.
	We have
	\[(\mathrm df)_{(x,y,z)} = \left(
	2(c-\sqrt{x^2+y^2})\cdot\frac{-2x}{2\sqrt{x^2+y^2}},
	2(c-\sqrt{x^2+y^2})\cdot\frac{-2y}{2\sqrt{x^2+y^2}},
	2z\right),\]
	and to show that this is a surjective linear map $\RR^3\to\RR$,
	we just need to verify that the components are not all zero
	when $f(x,y,z) = a^2$; for this, we check some cases.
	\begin{itemize}
		\ii If $z\ne 0$ then obviously $2z\ne 0$ as well,
		so $(\mathrm df)$ has rank $1$ as desired.
		\ii Otherwise suppose $z =0$. Then $(c-\sqrt{x^2+y^2})^2 = a^2$
		implies that $\sqrt{x^2+y^2}$ is either equal to $c-a$ or $c+a$.
		Since $c+a>c-a>0$, $x^2+y^2$ is nonzero, which implies that
		at least one of $x$ and $y$ is nonzero, say WLOG $x\ne0$.
		In this case the corresponding component of $(\mathrm df)$
		\[2(c-\sqrt{x^2+y^2})\cdot\frac{-x}{\sqrt{x^2+y^2}}\]
		is nonzero (noting that $\sqrt{x^2+y^2}$ is not equal to $c$
		because $a\ne0$) so again the derivative has rank $1$.
	\end{itemize}
	Therefore we can conclude that $f\inv(a)$ is a $C^\infty$
	submanifold of $\RR^3$ with dimension $3-1=2$.
	
	This set is a torus obtained by
	rotating a circle $\gamma$ about the origin,
	where $\gamma$ lies in the $xz$-plane such that
	its center is a distance of $c$ from the origin
	and its radius is $a$. To verify this, one can see that
	$f$ is invariant under rotation about the $z$-axis
	because such rotations preserve
	the distance $\sqrt{x^2+y^2}$ to the origin,
	and then check that the cross section in the $xz$-plane
	is indeed the described circle $(c-x)^2 + z^2 = a^2$
	(well, really it's the two circles $(c-|x|^2)+z^2=a^2$,
	but we can restrict to $x>0$ because the other circle
	$(c+x)^2+z^2=a^2$ is the one opposite this circle on the torus).
\end{soln}

\begin{problem}
	Let $f\colon\RR^4\to\RR^2$ be the function
	\[f(x,y,z,w) = (x^2-y^2+z^2-w^2, 2xy+2zw).\]
	Show that $f\inv(a,b)$ is a $2$-dimensional $C^\infty$
	submanifold of $\RR^4$ if $(a,b)\ne(0,0)$.
\end{problem}

\begin{soln}
	Clearly $f$ is $C^\infty$, because its components are polynomials
	which are $C^\infty$. To show that the pre-image of a nonzero point
	is a $2$-dimensional $C^\infty$ submanifold of $\RR^4$
	for each $x\in f\inv(a,b)$, we need to check that $(\mathrm df)_{(x,y,z,w)}$
	has rank $2$ for each $(x,y,z,w)\in f\inv(a,b)$. We have
	\[(\mathrm df)_{(x,y,z,w)} = \begin{pmatrix}
		2x&-2y&2z&-2w\\
		2y&2x&2w&2z.
	\end{pmatrix}\]
	If $(x,y)\ne(0,0)$ then the left two rows of the matrix are independent,
	and if $(z,w)\ne(0,0)$ then the right two rows are independent,
	so either of these forces $(\mathrm df)_{(x,y,z,w)}$ to be surjective.
	In particular the only way the linear map isn't surjective
	is if $(0,0,0,0)\in f\inv(a,b)$, but then this means
	$(a,b) = f(0,0,0,0) = (0,0)$, which we have forbidden.
	Therefore the derivative of $f$ is always a surjective map
	$\RR^4\to\RR^2$, so $f\inv(a,b)$ is a $C^\infty$ submanifold of $\RR^4$
	with dimension $4-2=2$ as desired.
\end{soln}

\begin{problem}
	Let $C\subseteq\RR^2$ be the \vocab{nodal curve}
	defined by the equation $x^2=y^2$. Is $C$ a $C^\infty$ submanifold
	of $\RR^2$?
\end{problem}

\begin{soln}
	No. We show that no neighborhood of $(0,0)$ in the subspace $C$
	is homeomorphic to $\RR^n$ for any positive integer $n$
	(and of course $C^\infty$ manifolds are also $C^0$ manifolds).
	Suppose otherwise: then, for some open $U\subseteq C$,
	there is a homeomorphism $f\colon C\to\RR^n$.
	Since $C$ is the union of the lines $y=x$ and $y=-x$,
	the set $C\setdiff\{(0,0)\}$ has four connected components:
	$\{(x,x) : x > 0\}$, $\{(x,x) : x < 0\}$, $\{(x, -x) : x>0\}$,
	and $\{(x,-x) : x < 0\}$.
	Then, since connectedness and disconnectedness
	is preserved by homemorphism, we get that $\RR^n\setdiff\{f(0)\}$
	must also have four connected components.
	However, this is impossible:
	\begin{itemize}
		\ii If $n = 1$, then deleting $f(0)$ from $\RR$
		gives two connected components, $(-\infty,f(0))$ and $(f(0),\infty)$.
		\ii If $n\ge 2$, then deleting $f(0)$ does not affect the
		connectedness of $\RR^n$.
	\end{itemize}
	This implies that $C$ does not have a $C^0$
	$n$-dimensional atlas for any $n$,
	so it is not a $C^0$ submanifold of $\RR^2$, let alone a $C^\infty$ one.
\end{soln}
















\end{document}